\section{总结与展望}

% 本章通常 1500–2500 字，总结全文并展望未来。

\subsection{研究结论}

% 用 3–5 条要点概括全文的核心发现。
% 每条结论对应第 4 章的一个主要结果。
% 不要引入新的数据或引用。

\subsection{研究贡献}

% 说明本研究在理论和实践两个层面的贡献。
% 与第 1 章"研究意义"呼应，但这里用完成时态。

\subsection{研究局限}

% 诚实说明研究的不足之处。
% 例如：样本范围有限、方法存在假设约束、数据时间跨度不够等。

\subsection{未来展望}

% 基于局限性提出后续研究方向。
% 可以指出值得深入探索的问题或可以改进的方法。
